% i18n_qa_mixed_fr_en_paper.tex
% French thesis with English abstract
% Expected detect_language result: "fr" (last babel option)
% Deliberate errors: missing NBSP before punctuation, wrong quote style
\documentclass[12pt,a4paper]{report}
\usepackage[utf8]{inputenc}
\usepackage[T1]{fontenc}
\usepackage[english,french]{babel}
\usepackage{amsmath,amssymb,amsthm}
\usepackage{graphicx}
\usepackage{booktabs}
\usepackage{hyperref}
\usepackage{algorithm}
\usepackage{algpseudocode}
\usepackage{tikz}
\usepackage{pgfplots}
\pgfplotsset{compat=1.18}

\newtheorem{theoreme}{Th\'eor\`eme}[chapter]
\newtheorem{lemme}[theoreme]{Lemme}
\newtheorem{proposition}[theoreme]{Proposition}
\newtheorem{corollaire}[theoreme]{Corollaire}
\newtheorem{definition}{D\'efinition}[chapter]

\title{Optimisation convexe sous contraintes stochastiques :\\
       Algorithmes adaptatifs et garanties de convergence}
\author{Marie-Claire Dubois\\
  \textit{Laboratoire de Math\'ematiques Appliqu\'ees}\\
  \textit{Universit\'e Paris-Saclay}}
\date{Janvier 2026}

\begin{document}
\maketitle

\begin{abstract}
\selectlanguage{english}
This thesis investigates adaptive algorithms for convex optimization under
stochastic constraints. We propose a novel framework that unifies proximal
gradient methods with variance reduction techniques, achieving optimal
convergence rates for both smooth and non-smooth objectives. Our main
contribution is a family of algorithms, termed Adaptive Stochastic Proximal
(ASP), that automatically adjust step sizes based on local curvature
estimates without requiring knowledge of the global Lipschitz constant.
We establish $O(1/\sqrt{T})$ convergence for the general convex case and
linear convergence under strong convexity assumptions. Numerical experiments
on large-scale machine learning problems, including LASSO regression and
support vector machines, demonstrate that ASP consistently outperforms
existing methods such as SGD with momentum, Adam, and AdaGrad.
Furthermore, we extend our analysis to the distributed setting, proving
that communication-efficient variants of ASP retain their theoretical
guarantees while reducing bandwidth requirements by a factor of $O(d)$
where $d$ is the problem dimension.

\medskip
\noindent\textbf{Keywords:} convex optimization, stochastic gradient descent,
proximal methods, variance reduction, adaptive step sizes, distributed optimization.
\selectlanguage{french}
\end{abstract}

\tableofcontents

\chapter{Introduction}

% ERROR: missing NBSP before colon (French typography requires NBSP before :;!?)
L'optimisation convexe constitue un pilier fondamental des math\'ematiques
appliqu\'ees modernes. Depuis les travaux fondateurs de Nesterov dans les
ann\'ees 1980, la th\'eorie de la complexit\'e des algorithmes d'optimisation
a connu des avanc\'ees consid\'erables. Dans cette th\`ese, nous nous
int\'eressons particuli\`erement aux probl\`emes d'optimisation sous
contraintes stochastiques: un cadre qui apparait naturellement dans de
nombreuses applications pratiques.

% ERROR: straight quotes instead of guillemets
Le probl\`eme g\'en\'eral que nous consid\'erons s'\'ecrit sous la forme
"standard" suivante:
\begin{equation}\label{eq:prob-principal}
  \min_{x \in \mathcal{X}} \left\{ F(x) \triangleq
    \mathbb{E}_{\xi \sim \mathcal{D}} \left[ f(x; \xi) \right]
    + \lambda \, g(x) \right\},
\end{equation}
o\`u $\mathcal{X} \subseteq \mathbb{R}^d$ est un ensemble convexe ferm\'e,
$f(\cdot; \xi)$ est une fonction convexe pour chaque r\'ealisation $\xi$
de la variable al\'eatoire, $g : \mathbb{R}^d \to \mathbb{R} \cup \{+\infty\}$
est une fonction convexe semi-continue inf\'erieurement, et $\lambda > 0$
est un param\`etre de r\'egularisation.

% ERROR: missing NBSP before semicolon
La motivation principale de ce travail provient des applications en
apprentissage automatique; en effet, la plupart des probl\`emes
d'entra\^inement de mod\`eles statistiques se formulent naturellement
sous la forme~\eqref{eq:prob-principal}. Par exemple, la r\'egression
logistique r\'egularis\'ee correspond au choix $f(x; \xi_i) =
\log(1 + \exp(-y_i \langle a_i, x \rangle))$ et $g(x) = \|x\|_1$.

Les contributions principales de cette th\`ese sont les suivantes:
\begin{enumerate}
  \item Nous proposons une famille d'algorithmes adaptatifs, appel\'ee
    ASP (Adaptive Stochastic Proximal), qui ajuste automatiquement les
    pas de descente en fonction de la courbure locale du probl\`eme.
  \item Nous \'etablissons des garanties de convergence optimales pour
    ASP dans le cas convexe g\'en\'eral ($O(1/\sqrt{T})$) et dans le
    cas fortement convexe (convergence lin\'eaire).
  \item Nous \'etendons notre analyse au cadre distribu\'e et montrons
    que les variantes efficaces en communication de ASP conservent leurs
    garanties th\'eoriques.
  \item Nous validons nos r\'esultats par des exp\'eriences num\'eriques
    approfondies sur des jeux de donn\'ees r\'eels de grande taille.
\end{enumerate}


\chapter{Pr\'eliminaires et \'etat de l'art}

\section{Cadre th\'eorique}

Nous commen\c{c}ons par rappeler les d\'efinitions et propri\'et\'es
fondamentales de l'optimisation convexe qui seront utilis\'ees tout
au long de cette th\`ese.

\begin{definition}[Convexit\'e forte]
Une fonction $f : \mathbb{R}^d \to \mathbb{R}$ est dite
$\mu$-fortement convexe (avec $\mu > 0$) si pour tout
$x, y \in \mathbb{R}^d$ et tout $\alpha \in [0, 1]$:
\begin{equation}\label{eq:strong-cvx}
  f(\alpha x + (1-\alpha) y) \leq \alpha f(x) + (1-\alpha) f(y)
    - \frac{\mu}{2} \alpha (1-\alpha) \|x - y\|^2.
\end{equation}
\end{definition}

\begin{definition}[Lissit\'e]
Une fonction diff\'erentiable $f : \mathbb{R}^d \to \mathbb{R}$ est
dite $L$-lisse si son gradient est $L$-lipschitzien:
\begin{equation}\label{eq:lipschitz-grad}
  \|\nabla f(x) - \nabla f(y)\| \leq L \|x - y\|, \quad
  \forall x, y \in \mathbb{R}^d.
\end{equation}
\end{definition}

Le nombre de condition $\kappa = L / \mu$ joue un r\^ole central dans
l'analyse de la vitesse de convergence des algorithmes d'optimisation.
Plus $\kappa$ est grand, plus le probl\`eme est mal conditionn\'e et
plus la convergence est lente.

\section{M\'ethodes de gradient stochastique}

% ERROR: missing NBSP before exclamation mark
La descente de gradient stochastique (SGD), introduite par Robbins et
Monro en 1951, est l'algorithme le plus fondamental pour r\'esoudre le
probl\`eme~\eqref{eq:prob-principal}. L'id\'ee est remarquablement simple!
\`A chaque it\'eration $t$, on tire un \'echantillon $\xi_t$ et on
effectue la mise \`a jour:
\begin{equation}\label{eq:sgd}
  x_{t+1} = \Pi_{\mathcal{X}} \left( x_t - \eta_t \nabla f(x_t; \xi_t) \right),
\end{equation}
o\`u $\Pi_{\mathcal{X}}$ d\'esigne la projection sur $\mathcal{X}$ et
$\eta_t > 0$ est le pas de descente.

\begin{theoreme}[Convergence de SGD, \cite{nemirovski2009robust}]
Sous les hypoth\`eses standard (convexit\'e de $f$, gradient born\'e
$\mathbb{E}[\|\nabla f(x; \xi)\|^2] \leq G^2$), le choix
$\eta_t = c / \sqrt{t}$ assure:
\begin{equation}
  \mathbb{E}[F(\bar{x}_T)] - F(x^*) \leq \frac{C}{\sqrt{T}},
\end{equation}
o\`u $\bar{x}_T = \frac{1}{T} \sum_{t=1}^T x_t$ et $C$ d\'epend de
$G$, $c$ et $\|x_0 - x^*\|$.
\end{theoreme}

La limitation principale de SGD r\'eside dans la variance du gradient
stochastique, qui impose un pas de descente d\'ecroissant et conduit
\`a une convergence sous-lin\'eaire m\^eme pour les probl\`emes
fortement convexes.

\section{R\'eduction de variance}

Les m\'ethodes de r\'eduction de variance, telles que SVRG
\cite{johnson2013accelerating} et SAGA \cite{defazio2014saga},
permettent de surmonter cette limitation en construisant des estimateurs
du gradient \`a variance d\'ecroissante.

\begin{lemme}[R\'eduction de variance SVRG]
L'estimateur SVRG d\'efini par
$v_t = \nabla f(x_t; \xi_t) - \nabla f(\tilde{x}; \xi_t) + \nabla F(\tilde{x})$
satisfait:
\begin{equation}
  \mathbb{E}[\|v_t - \nabla F(x_t)\|^2] \leq 4L \left( F(x_t) - F(x^*)
    + F(\tilde{x}) - F(x^*) \right).
\end{equation}
\end{lemme}

Cette propri\'et\'e permet d'obtenir une convergence lin\'eaire pour
les probl\`emes fortement convexes avec un pas de descente constant,
ce qui repr\'esente une am\'elioration significative par rapport \`a SGD.


\chapter{Algorithmes adaptatifs proximaux stochastiques}

\section{Formulation g\'en\'erale}

Nous introduisons maintenant notre famille d'algorithmes ASP. L'id\'ee
principale est de combiner les avantages des m\'ethodes proximales avec
une estimation adaptative de la courbure locale.

\begin{algorithm}
\caption{ASP -- Adaptive Stochastic Proximal}\label{alg:asp}
\begin{algorithmic}[1]
\Require Point initial $x_0$, param\`etres $\beta_1, \beta_2 \in (0,1)$
\State Initialiser $m_0 = 0$, $v_0 = \epsilon \mathbf{I}_d$
\For{$t = 0, 1, 2, \ldots, T-1$}
  \State Tirer $\xi_t \sim \mathcal{D}$
  \State $g_t = \nabla f(x_t; \xi_t)$
  \State $m_t = \beta_1 m_{t-1} + (1-\beta_1) g_t$
  \State $v_t = \beta_2 v_{t-1} + (1-\beta_2) g_t \odot g_t$
  \State $\hat{m}_t = m_t / (1 - \beta_1^{t+1})$
  \State $\hat{v}_t = v_t / (1 - \beta_2^{t+1})$
  \State $H_t = \mathrm{diag}(\sqrt{\hat{v}_t} + \epsilon)$
  \State $x_{t+1} = \mathrm{prox}_{g}^{H_t} \left( x_t - H_t^{-1} \hat{m}_t \right)$
\EndFor
\State \Return $\bar{x}_T = \frac{1}{T} \sum_{t=0}^{T-1} x_t$
\end{algorithmic}
\end{algorithm}

% ERROR: missing NBSP before question mark
L'op\'erateur proximal g\'en\'eralis\'e utilis\'e \`a la ligne~10 est
d\'efini par:
\begin{equation}\label{eq:prox-gen}
  \mathrm{prox}_{g}^{H}(y) = \arg\min_{x \in \mathcal{X}} \left\{
    g(x) + \frac{1}{2} (x - y)^\top H (x - y) \right\}.
\end{equation}
Pourquoi cette formulation est-elle avantageuse? Elle combine la
r\'egularisation implicite de l'op\'erateur proximal avec l'adaptation
de la m\'etrique \`a la g\'eom\'etrie locale du probl\`eme.

\section{Analyse de convergence}

\begin{theoreme}[Convergence de ASP -- cas convexe]\label{thm:asp-convex}
Soit $F$ une fonction convexe et $L$-lisse. Sous les hypoth\`eses~(H1)-(H3)
d\'efinies dans l'Annexe~A, l'algorithme ASP avec $\beta_1 = 1 - 1/\sqrt{T}$
et $\beta_2 = 1 - 1/T$ satisfait:
\begin{equation}
  \mathbb{E}[F(\bar{x}_T)] - F(x^*) \leq \frac{C_1 \sqrt{d}}{\sqrt{T}}
    \left( 1 + \log T \right),
\end{equation}
o\`u $C_1$ d\'epend de $L$, $\|x_0 - x^*\|$ et des moments du bruit.
\end{theoreme}

\begin{proof}[D\'emonstration (esquisse)]
La preuve repose sur trois ingr\'edients principaux. Premi\`erement, nous
\'etablissons un lemme de descente adapt\'e \`a la m\'etrique variable $H_t$.
Deuxi\`emement, nous bornons l'erreur d'approximation due \`a la correction
de biais. Troisi\`emement, nous utilisons une analyse de regret \`a la
mani\`ere de Zinkevich pour conclure.

Soit $r_t = F(x_t) - F(x^*)$. Par convexit\'e de $F$ et la propri\'et\'e
de l'op\'erateur proximal:
\begin{align}
  r_{t+1} &\leq \langle \nabla F(x_t), x_t - x^* \rangle
    - \frac{1}{2} \|x_{t+1} - x_t\|_{H_t}^2 \nonumber \\
  &\leq r_t - \frac{\mu_{\min}(H_t)}{2} \|x_{t+1} - x_t\|^2
    + \langle g_t - \nabla F(x_t), x_t - x^* \rangle.
\end{align}
En sommant de $t=0$ \`a $T-1$ et en prenant l'esp\'erance, on obtient le
r\'esultat apr\`es des majorations techniques d\'etaill\'ees dans
l'Annexe~A.
\end{proof}

\begin{theoreme}[Convergence lin\'eaire -- cas fortement convexe]\label{thm:asp-strong}
Si en outre $F$ est $\mu$-fortement convexe avec $\kappa = L/\mu$, alors
ASP avec red\'emarrage p\'eriodique (p\'eriode $T_0 = O(\kappa \sqrt{d})$)
satisfait apr\`es $K$ \'epoques:
\begin{equation}
  \mathbb{E}[F(x^{(K)})] - F(x^*) \leq
    \left( 1 - \frac{c}{\kappa \sqrt{d}} \right)^K
    \left( F(x_0) - F(x^*) \right).
\end{equation}
\end{theoreme}


\chapter{Extension au cadre distribu\'e}

\section{Mod\`ele de communication}

Nous consid\'erons un r\'eseau de $n$ agents, chacun ayant acc\`es \`a
une distribution locale $\mathcal{D}_i$. Le probl\`eme global s'\'ecrit:
\begin{equation}\label{eq:dist-prob}
  \min_{x \in \mathcal{X}} \frac{1}{n} \sum_{i=1}^n
    \mathbb{E}_{\xi \sim \mathcal{D}_i} [f_i(x; \xi)] + \lambda g(x).
\end{equation}

% ERROR: missing NBSP before colon
La contrainte principale dans ce cadre est la communication entre agents:
chaque \'echange de vecteurs dans $\mathbb{R}^d$ a un co\^ut proportionnel
\`a $d$, ce qui peut devenir prohibitif pour les probl\`emes de grande
dimension.

\section{Compression par quantification}

Nous proposons une variante de ASP avec compression des gradients
communiqu\'es. Soit $Q : \mathbb{R}^d \to \mathbb{R}^d$ un op\'erateur
de quantification satisfaisant:
\begin{equation}
  \mathbb{E}[\|Q(x) - x\|^2] \leq (1-\delta) \|x\|^2,
\end{equation}
pour un certain $\delta \in (0, 1]$.\footnote{Cette propri\'et\'e est
v\'erifi\'ee par de nombreux op\'erateurs classiques, notamment la
quantification stochastique \`a $b$ bits avec $\delta = \min(b/d, 1)$.}

\begin{proposition}[Convergence de ASP distribu\'e compress\'e]
L'algorithme ASP distribu\'e avec compression $Q$ converge \`a la vitesse:
\begin{equation}
  \mathbb{E}[F(\bar{x}_T)] - F(x^*) \leq
    \frac{C_2}{\delta} \cdot \frac{\sqrt{d}}{\sqrt{nT}},
\end{equation}
ce qui correspond \`a une acc\'el\'eration lin\'eaire en $n$ (speed-up
lin\'eaire).
\end{proposition}

Le tableau~\ref{tab:comparison} r\'esume les r\'esultats de convergence
pour les diff\'erentes variantes de notre algorithme.

\begin{table}[htbp]
\centering
\caption{Comparaison des vitesses de convergence}\label{tab:comparison}
\begin{tabular}{@{}lccc@{}}
\toprule
Algorithme & Cas convexe & Cas fort. convexe & Communication \\
\midrule
SGD        & $O(1/\sqrt{T})$ & $O(1/T)$ & $O(dT)$ \\
Adam       & $O(1/\sqrt{T})$ & $O(\sqrt{d}/T)$ & --- \\
ASP        & $O(\sqrt{d}/\sqrt{T})$ & $O(\exp(-T/\kappa\sqrt{d}))$ & --- \\
ASP-Dist   & $O(\sqrt{d}/\sqrt{nT})$ & $O(\exp(-nT/\kappa\sqrt{d}))$ & $O(dT)$ \\
ASP-Comp   & $O(\sqrt{d}/\delta\sqrt{nT})$ & $O(\exp(-\delta nT/\kappa\sqrt{d}))$ & $O(bT)$ \\
\bottomrule
\end{tabular}
\end{table}


\chapter{Exp\'eriences num\'eriques}

\section{Configuration exp\'erimentale}

% ERROR: missing NBSP before semicolon
Nous \'evaluons nos algorithmes sur trois probl\`emes classiques
d'apprentissage automatique; la r\'egression LASSO, les machines \`a
vecteurs de support (SVM) et la r\'egression logistique r\'egularis\'ee.

Les jeux de donn\'ees utilis\'es sont d\'ecrits dans le
tableau~\ref{tab:datasets}.

\begin{table}[htbp]
\centering
\caption{Jeux de donn\'ees utilis\'es dans les exp\'eriences}\label{tab:datasets}
\begin{tabular}{@{}lrrr@{}}
\toprule
Jeu de donn\'ees & $n$ (\'echantillons) & $d$ (dimension) & Taille (Mo) \\
\midrule
rcv1-train  & 20\,242  & 47\,236  & 87 \\
news20      & 15\,935  & 62\,061  & 120 \\
covtype     & 581\,012 & 54       & 72 \\
MNIST-8M    & 8\,100\,000 & 784   & 24\,000 \\
\bottomrule
\end{tabular}
\end{table}

\section{R\'esultats sur la r\'egression LASSO}

Le probl\`eme LASSO consiste \`a r\'esoudre:
\begin{equation}
  \min_{x \in \mathbb{R}^d} \frac{1}{2n} \|Ax - b\|_2^2 + \lambda \|x\|_1,
\end{equation}
o\`u $A \in \mathbb{R}^{n \times d}$ et $b \in \mathbb{R}^n$ sont les
donn\'ees d'entra\^inement. Nous fixons $\lambda = 10^{-4}$ et comparons
ASP avec SGD, Adam, AdaGrad et SVRG.

La figure~\ref{fig:lasso-convergence} illustre la convergence de l'erreur
relative $[F(x_t) - F(x^*)] / F(x^*)$ en fonction du nombre d'\'epoques.

\begin{figure}[htbp]
\centering
\begin{tikzpicture}
\begin{semilogyaxis}[
  width=0.85\textwidth,
  height=6cm,
  xlabel={\'Epoques},
  ylabel={Erreur relative},
  legend pos=north east,
  grid=major,
  xmin=0, xmax=100,
]
\addplot[blue, thick] coordinates {
  (0,1) (10,0.3) (20,0.08) (30,0.02) (40,0.005) (50,0.001)
  (60,3e-4) (70,8e-5) (80,2e-5) (90,5e-6) (100,1e-6)
};
\addplot[red, thick, dashed] coordinates {
  (0,1) (10,0.5) (20,0.2) (30,0.1) (40,0.05) (50,0.03)
  (60,0.015) (70,0.008) (80,0.005) (90,0.003) (100,0.002)
};
\addplot[green!60!black, thick, dotted] coordinates {
  (0,1) (10,0.4) (20,0.15) (30,0.06) (40,0.02) (50,0.008)
  (60,0.003) (70,0.001) (80,4e-4) (90,1.5e-4) (100,5e-5)
};
\legend{ASP, SGD+Momentum, SVRG}
\end{semilogyaxis}
\end{tikzpicture}
\caption{Convergence de l'erreur relative sur rcv1-train (LASSO,
  $\lambda = 10^{-4}$).}\label{fig:lasso-convergence}
\end{figure}

\section{R\'esultats distribu\'es}

Nous simulons un r\'eseau de $n = 16$ agents et comparons ASP-Dist
et ASP-Comp (avec quantification \`a 4 bits) avec des m\'ethodes
distribu\'ees de r\'ef\'erence.

\begin{table}[htbp]
\centering
\caption{Temps d'entra\^inement (secondes) pour atteindre $\epsilon = 10^{-4}$
  sur covtype ($n=16$ agents)}\label{tab:dist-results}
\begin{tabular}{@{}lrrr@{}}
\toprule
M\'ethode & Temps (s) & Communication (Go) & Speed-up \\
\midrule
ASP (centralis\'e)  & 42.3  & ---  & 1.0$\times$ \\
ASP-Dist            & 3.8   & 1.2  & 11.1$\times$ \\
ASP-Comp (4 bits)   & 4.2   & 0.15 & 10.1$\times$ \\
D-PSGD              & 5.1   & 1.2  & 8.3$\times$ \\
CHOCO-SGD           & 4.7   & 0.18 & 9.0$\times$ \\
\bottomrule
\end{tabular}
\end{table}


\chapter{Conclusion et perspectives}

% ERROR: straight quotes instead of guillemets
Nous avons pr\'esent\'e dans cette th\`ese la famille d'algorithmes ASP,
qui offre un cadre unifi\'e pour l'optimisation convexe stochastique
adaptative. Nos r\'esultats th\'eoriques montrent que ASP atteint des
vitesses de convergence "optimales" au sens minimax, tout en s'adaptant
automatiquement \`a la g\'eom\'etrie locale du probl\`eme.

Les perspectives de ce travail sont nombreuses. Premi\`erement, l'extension
aux probl\`emes non convexes, qui apparaissent fr\'equemment dans
l'entra\^inement des r\'eseaux de neurones profonds, constitue un d\'efi
majeur. Deuxi\`emement, l'int\'egration de techniques de pr\'econditionnement
de second ordre pourrait am\'eliorer la convergence dans les r\'egimes
mal conditionn\'es. Troisi\`emement, l'application \`a des probl\`emes
d'optimisation sous contraintes plus g\'en\'erales, incluant des contraintes
d'in\'egalit\'e non lin\'eaires, reste \`a explorer.

\section*{Remerciements}

% Footnote in French
L'auteure remercie son directeur de th\`ese, le Professeur Jean-Luc
Martin\footnote{Actuellement en d\'el\'egation au CNRS, UMR 7641.},
pour son encadrement exceptionnel et ses conseils avis\'es tout au long
de ce travail. Elle remercie \'egalement les membres du jury pour leurs
remarques constructives lors de la soutenance.


\begin{thebibliography}{99}

\bibitem{nemirovski2009robust}
A.~Nemirovski, A.~Juditsky, G.~Lan, and A.~Shapiro.
\newblock Robust stochastic approximation approach to stochastic programming.
\newblock \textit{SIAM Journal on Optimization}, 19(4):1574--1609, 2009.

\bibitem{johnson2013accelerating}
R.~Johnson and T.~Zhang.
\newblock Accelerating stochastic gradient descent using predictive variance
  reduction.
\newblock In \textit{Advances in Neural Information Processing Systems (NeurIPS)},
  pages 315--323, 2013.

\bibitem{defazio2014saga}
A.~Defazio, F.~Bach, and S.~Lacoste-Julien.
\newblock {SAGA}: A fast incremental gradient method with support for
  non-strongly convex composite objectives.
\newblock In \textit{Advances in Neural Information Processing Systems (NeurIPS)},
  pages 1646--1654, 2014.

\bibitem{kingma2015adam}
D.~P. Kingma and J.~Ba.
\newblock Adam: A method for stochastic optimization.
\newblock In \textit{International Conference on Learning Representations
  (ICLR)}, 2015.

\bibitem{koloskova2019decentralized}
A.~Koloskova, S.~Stich, and M.~Jaggi.
\newblock Decentralized stochastic optimization and gossip algorithms with
  compressed communication.
\newblock In \textit{International Conference on Machine Learning (ICML)},
  pages 3478--3487, 2019.

\bibitem{duchi2011adaptive}
J.~Duchi, E.~Hazan, and Y.~Singer.
\newblock Adaptive subgradient methods for online learning and stochastic
  optimization.
\newblock \textit{Journal of Machine Learning Research}, 12:2121--2159, 2011.

\bibitem{nesterov2004introductory}
Y.~Nesterov.
\newblock \textit{Introductory Lectures on Convex Optimization: A Basic Course}.
\newblock Springer, 2004.

\bibitem{bottou2018optimization}
L.~Bottou, F.~E. Curtis, and J.~Nocedal.
\newblock Optimization methods for large-scale machine learning.
\newblock \textit{SIAM Review}, 60(2):223--311, 2018.

\end{thebibliography}

\end{document}
